\section{Why one finite nilpotent quotient is insufficient}
\label{sec:peripheral-finite-obstruction}

We exhibit a finite group generated by a product-one triple and a normal
subgroup whose nilpotent residual does not suffice to correct prescribed
conjugator cosets. The quotient automorphism in the example is the identity,
and the conjugated first two elements still generate the whole group.
The obstruction is an explicit central sign in a central product of
$\operatorname{SL}_2(\F_5)$ and $\operatorname{SL}_2(\F_3)$. A quaternion quotient detects the same
obstruction, explaining why the example cannot arise from the corresponding
full pronilpotent quotient of a free profinite group.

For a group $H$, write $\gamma_{\infty}(H)=\bigcap_{n\geq1}\gamma_n(H)$, with
closed lower central terms when $H$ is profinite. If $D\leq G$, write
$g^D=\{dgd^{-1}:d\in D\}$.

\begin{theorem}\label{thm:peripheral-finite-obstruction}
There exist a finite group $G=\langle g_1,g_2\rangle$, a normal subgroup
$N\trianglelefteq G$, and elements $a_1,a_2,a_3\in G$ with the following
properties. Put $g_3=(g_1g_2)^{-1}$, $D=\gamma_{\infty}(N)$, and
$h_i=a_ig_ia_i^{-1}$. Then
\[
|G|=1440,\qquad |N|=480,\qquad |D|=120,\qquad G/D\cong A_4,
\]
and
\[
h_iD=g_iD\quad(1\leq i\leq3),\qquad
G=\langle h_1,h_2\rangle,\qquad
1\notin h_1^D h_2^D h_3^D.
\]
One may take $a_2=a_3=1$. In particular, the finite quotient datum is an
identity automorphism with specified peripheral conjugators, but those
conjugator cosets cannot be corrected by $D$.
\end{theorem}

\begin{lemma}\label{lem:peripheral-central-sign}
In $S=\operatorname{SL}_2(\F_5)$, put
\[
A=\begin{pmatrix}0&1\\4&1\end{pmatrix},\qquad
B=\begin{pmatrix}0&1\\4&3\end{pmatrix},\qquad
C=\begin{pmatrix}2&2\\1&4\end{pmatrix}.
\]
Then $S=\langle A,B\rangle$, $ABC=I$, and
\[
I\in A^S B^S C^S,\qquad -I\notin A^S B^S C^S.
\]
\end{lemma}

\begin{proof}
We first establish generation, then exclude the opposite central sign.
Direct multiplication gives $C=(AB)^{-1}$ and
\[
(A^{-1}B)^2=\begin{pmatrix}1&1\\0&1\end{pmatrix},\qquad
(AB^{-1})^2=\begin{pmatrix}1&0\\1&1\end{pmatrix}.
\]
The elementary matrices on the right generate $\operatorname{SL}_2(\F_5)$ by row
operations, proving generation and the first inclusion.

Suppose $uvw=-I$, with $u\in A^S$, $v\in B^S$, and $w\in C^S$.
Simultaneous conjugation permits $u=A$. Since
$\operatorname{tr} A=\operatorname{tr} C=1$ and $\operatorname{tr} B=-2$, the equation forces
$\operatorname{tr}(Av)=\operatorname{tr}(-w^{-1})=-1$. Write
\[
v=-I+\begin{pmatrix}r&s\\t&-r\end{pmatrix}.
\]
The determinant condition gives $r^2+st=0$, while the trace condition
gives $t-s-r=0$. Hence $r^2+rs+s^2=0$. This quadratic form has no
nonzero zero over $\F_5$, because its discriminant $-3=2$ is a
nonsquare. Thus $r=s=t=0$ and $v=-I$, contradicting $v\in B^S$.
\end{proof}

\begin{proof}[Proof of Theorem~\ref{thm:peripheral-finite-obstruction}]
We place the sign obstruction in a central product whose nilpotent
residual is the first factor. Let $S,A,B,C$ be as in Lemma~\ref{lem:peripheral-central-sign}. The group $S$ is perfect:
conjugation by $\operatorname{diag}(2,3)$ multiplies the parameter in
each elementary unipotent subgroup by $4$, so the elementary matrix
of parameter $1$ is the commutator with the elementary matrix of
parameter $2$. Since these elementary matrices generate $S$, we have
$[S,S]=S$.

Over $\F_3$, put
\[
E=\operatorname{SL}_2(\F_3),\qquad
i=\begin{pmatrix}0&1\\2&0\end{pmatrix},\qquad
t=\begin{pmatrix}1&1\\0&1\end{pmatrix},\qquad
j=\begin{pmatrix}1&1\\1&2\end{pmatrix}.
\]
The matrices satisfy $i^2=j^2=-I$ and $ji=-ij$. The subgroup
\[
Q=\{\pm I,\pm i,\pm j,\pm ij\}
\]
is quaternion of order eight and is preserved by conjugation by $t$.
Moreover $i,t$ generate $E$: $t$ is an elementary upper unipotent, and
conjugation by $i$ supplies a lower unipotent. Thus $Q$ is normal in
$E$, and $E/Q$ has order three. In particular, $E$ is solvable. The
quotients $E/\{\pm I\}$ and $Q/\{\pm I\}$ are $A_4$ and its normal
Klein four subgroup.

Form
\[
G=(S\times E)/\langle(-I,-I)\rangle.
\]
Use brackets for its elements. Let $D$ be the image of $S\times\{I\}$
and $N$ the image of $S\times Q$. The embedding $S\to D$ is
injective. Since $|S|=120$ and $|E|=24$, the displayed group orders
follow. Also $G/D\cong A_4$ and $N/D\cong V_4$. The identities
$[S,S]=S$ and $[Q,Q]=\{\pm I\}$ give
\[
[N,N]=D,\qquad [D,N]=D,
\]
so $\gamma_{\infty}(N)=D$.

Define
\[
g_1=[(A,i)],\qquad g_2=[(B,t)],\qquad
g_3=[(C,(it)^{-1})].
\]
Their product is one. The subgroup of $S\times E$ generated by
$(A,i),(B,t)$ projects onto each factor. Any common quotient of the
two factors is both perfect and solvable, hence trivial. The elementary
subdirect-product argument therefore makes this subgroup all of
$S\times E$, proving $G=\langle g_1,g_2\rangle$.

Take $a_1=[(I,j)]$ and $a_2=a_3=1$. Then
\[
h_1=[(A,-i)]=[(-A,i)],\qquad h_2=g_2,\qquad h_3=g_3.
\]
The images of $i,j$ commute in $Q/\{\pm I\}$, so $h_iD=g_iD$ for
each $i$. Moreover $B^5=-I$, so $-A,B$ generate $S$; the same
subdirect-product argument proves $G=\langle h_1,h_2\rangle$.

Every element $d_k\in D$ has a unique representative $(s_k,I)$.
The product of the three resulting $D$-conjugates of the $h_k$ is
\[
[( -uvw,I)],\qquad u\in A^S,\quad v\in B^S,\quad w\in C^S.
\]
This is the identity in $G$ precisely when $uvw=-I$: its second
coordinate is $I$, so it cannot represent the nonidentity element
$(-I,-I)$ of the central subgroup being factored out. Lemma~\ref{lem:peripheral-central-sign}
excludes this possibility.
\end{proof}

The full pronilpotent quotient retains more information than the quotient
$G/D$ in Theorem~\ref{thm:peripheral-finite-obstruction}. The following observation makes the
distinction explicit.

\begin{proposition}\label{prop:peripheral-full-quotient-excludes}
Let $\Lambda=\widehat F_2=\langle x_1,x_2\rangle$, put
$x_3=(x_1x_2)^{-1}$, and let $q:\Lambda\twoheadrightarrow G$ send
$x_1,x_2$ to $g_1,g_2$. Set $U=q^{-1}(N)$, $R=\gamma_{\infty}(U)$, and
$\bar\Lambda=\Lambda/R$, and write $\bar x_i$ for the image of $x_i$.
There is no homomorphism
$\alpha:\bar\Lambda\to\bar\Lambda$ and conjugators $\bar c_i$ such
that
\[
\alpha(\bar x_i)=\bar c_i\bar x_i\bar c_i^{-1}
\]
and the images of the $\bar c_i$ in $G/D$ are $a_1D,1,1$.
\end{proposition}

\begin{proof}
Continuous maps of profinite groups commute with the nilpotent residual
when the target is finite: images commute with each lower central term,
and compactness commutes the image with their descending intersection.
Hence $q(R)=\gamma_{\infty}(N)=D$.

There is also a homomorphism
\[
e:\Lambda\longrightarrow E,\qquad e(x_1)=i,\quad e(x_2)=t.
\]
Its composite to $E/\{\pm I\}$ agrees with the composite of $q$ to
$G/D$. Therefore $e(U)=Q$, which is nilpotent, and $e(R)=1$. Write
$\bar e:\bar\Lambda\to E$ for the resulting homomorphism.

The asserted finite conjugator images would imply
\[
\bar e(\bar c_1)\in\{j,-j\},\qquad
\bar e(\bar c_2),\bar e(\bar c_3)\in\{I,-I\}.
\]
Consequently the images under $\bar e\alpha$ of the three peripheral
generators would be $-i$, $t$, and $(it)^{-1}$. Their product is $-I$,
contradicting $x_1x_2x_3=1$.
\end{proof}
